\documentclass[10pt, letterpaper]{article}

% ─── Packages ────────────────────────────────────────────────────────────────
\usepackage[top=1.1in, bottom=1.1in, left=1.25in, right=1.25in]{geometry}
\usepackage{amsmath, amssymb}
\usepackage{graphicx}
\usepackage[hidelinks, colorlinks=true, linkcolor=black, citecolor=black, urlcolor=blue]{hyperref}
\usepackage{booktabs}
\usepackage{array}
\usepackage{longtable}
\usepackage{enumitem}
\usepackage{float}
\usepackage{fancyhdr}
\usepackage{setspace}
\usepackage{hyperref}
\usepackage{xcolor}
\usepackage{listings}
\usepackage{caption}
\usepackage{cite}
\usepackage{url}
\usepackage{microtype}
\usepackage{parskip}
\usepackage{titlesec}
\usepackage{mdframed}
\usepackage{tocloft}
\usepackage{abstract}
\usepackage{lipsum}

% ─── Fonts ───────────────────────────────────────────────────────────────────
\usepackage[T1]{fontenc}
\usepackage{palatino}

% ─── Section Title Formatting ────────────────────────────────────────────────
\titleformat{\section}
  {\large\bfseries}{\thesection.}{0.6em}{}[\vspace{2pt}\titlerule\vspace{4pt}]

\titleformat{\subsection}
  {\normalsize\bfseries}{\thesubsection}{0.6em}{}

\titleformat{\subsubsection}
  {\normalsize\itshape}{\thesubsubsection}{0.6em}{}

\titlespacing*{\section}{0pt}{18pt}{8pt}
\titlespacing*{\subsection}{0pt}{12pt}{4pt}
\titlespacing*{\subsubsection}{0pt}{8pt}{3pt}

% ─── TOC Formatting ──────────────────────────────────────────────────────────
\renewcommand{\cftsecfont}{\bfseries}
\renewcommand{\cftsecpagefont}{\bfseries}
\setlength{\cftbeforesecskip}{4pt}

% ─── Header / Footer ─────────────────────────────────────────────────────────
\pagestyle{fancy}
\fancyhf{}
\renewcommand{\headrulewidth}{0.4pt}
\renewcommand{\footrulewidth}{0.4pt}
\lhead{\small\textit{MediQuery: A RAG System for Medicare Policy Q\&A}}
\rhead{\small\textit{GEN AI $\cdot$ MSBA 285}}
\cfoot{\thepage}

% ─── Code Style ──────────────────────────────────────────────────────────────
\definecolor{codebg}{RGB}{248,248,248}
\definecolor{codegreen}{RGB}{0,128,0}
\definecolor{codegray}{RGB}{128,128,128}
\definecolor{codeblue}{RGB}{0,0,180}
\definecolor{codeorange}{RGB}{180,80,0}

\lstdefinestyle{pythonstyle}{
    backgroundcolor=\color{codebg},
    commentstyle=\color{codegreen}\itshape,
    keywordstyle=\color{codeblue}\bfseries,
    stringstyle=\color{codeorange},
    basicstyle=\ttfamily\small,
    breaklines=true,
    captionpos=b,
    numbers=left,
    numbersep=8pt,
    numberstyle=\tiny\color{codegray},
    showstringspaces=false,
    tabsize=2,
    frame=single,
    framerule=0.5pt,
    rulecolor=\color{gray!50},
    xleftmargin=16pt,
    framexleftmargin=16pt,
    aboveskip=10pt,
    belowskip=10pt
}
\lstset{style=pythonstyle}

% ─── TODO Box ────────────────────────────────────────────────────────────────
\newmdenv[
  topline=true, bottomline=true, rightline=true, leftline=true,
  linewidth=1pt, linecolor=gray!60,
  backgroundcolor=gray!8,
  innertopmargin=8pt, innerbottommargin=8pt,
  innerleftmargin=10pt, innerrightmargin=10pt,
  skipabove=8pt, skipbelow=8pt
]{todobox}

% ─── Spacing ─────────────────────────────────────────────────────────────────
\setlength{\parskip}{6pt}
\setlength{\parindent}{0pt}
\onehalfspacing

% ─────────────────────────────────────────────────────────────────────────────
%  TITLE PAGE
% ─────────────────────────────────────────────────────────────────────────────
\begin{document}

\begin{titlepage}
    \centering
    \vspace*{2cm}

    {\large \textsc{University of Texas at Austin} \par}
    {\large \textsc{McCombs School of Business} \par}
    {\large GENERATIVE AI 285N $\cdot$ MSBA 285 \par}

    \vspace{1.5cm}
    \rule{\linewidth}{0.6pt}\\[0.4cm]

    {\LARGE \textbf{MediQuery} \par}
    \vspace{0.3cm}
    {\large A Retrieval-Augmented Generation System\\[4pt]
    for Medicare Policy Question Answering \par}

    \rule{\linewidth}{0.6pt}

    \vspace{2cm}

    {\normalsize
    \begin{tabular}{c}
        Mohar Chaudhuri \\[4pt]
        Shruthi Chembu Kuppuswamy\\[4pt]
        Sanchal Nachappa\\[4pt]
        Janani Vakkanti\\[4pt]
        Ryota Yokoyama
    \end{tabular}
    \par}

    \vfill

    {\normalsize March 15, 2026 \par}
    \vspace{0.5cm}
\end{titlepage}

% ─────────────────────────────────────────────────────────────────────────────
%  ABSTRACT
% ─────────────────────────────────────────────────────────────────────────────
\newpage
\thispagestyle{fancy}

\begin{center}
    {\large \textbf{Abstract}}
\end{center}

\vspace{0.3cm}

\noindent
Navigating Medicare coverage policy is a persistent challenge for healthcare providers,
administrators, and patients alike. Thousands of federal and state-level policy documents
govern what procedures and treatments are reimbursable under Medicare, yet these documents
are fragmented across government repositories, written in dense regulatory language, and
frequently updated. This project presents \textbf{MediQuery}, a Retrieval-Augmented Generation
(RAG) system that answers natural language questions about Medicare coverage by grounding
responses in real, official policy documents --- specifically, National Coverage Determinations
(NCDs) and Local Coverage Determinations (LCDs) sourced directly from the Centers for
Medicare \& Medicaid Services (CMS).

We describe the full system pipeline: raw document acquisition from CMS, preprocessing and
structured chunking of 1,162 policy documents into 3,377 semantically coherent chunks,
semantic embedding via \texttt{BAAI/bge-base-en-v1.5}, two-stage retrieval with cross-encoder
reranking via \texttt{BAAI/bge-reranker-v2-m3}, and grounded answer generation using
Mistral-7B-Instruct. The system is evaluated across three configurations --- a baseline LLM,
RAG, and RAG with query rewriting --- measuring retrieval recall, answer faithfulness,
citation accuracy, and hallucination rate. Our results demonstrate that RAG substantially
reduces hallucination and improves citation accuracy compared to an ungrounded language
model baseline.

\vspace{1cm}

\newpage

% ─────────────────────────────────────────────────────────────────────────────
%  TABLE OF CONTENTS
% ─────────────────────────────────────────────────────────────────────────────
\tableofcontents
 
\newpage
 
% ─────────────────────────────────────────────────────────────────────────────
\section{Introduction}
% ─────────────────────────────────────────────────────────────────────────────
 
Medicare is the United States federal health insurance program covering approximately
65 million Americans aged 65 and older or living with qualifying disabilities. Determining
what Medicare covers for a given patient, procedure, or geographic location requires
consulting a complex and constantly evolving set of policy documents. These documents
are maintained by the Centers for Medicare \& Medicaid Services (CMS) and exist in two
primary forms: \texti
t{National Coverage Determinations} (NCDs), which establish
federal-level coverage applicable uniformly across all states, and \textit{Local Coverage
Determinations} (LCDs), which define regional rules set by Medicare Administrative
Contractors (MACs) for specific jurisdictions.
 
In practice, a provider trying to answer a question such as \textit{``Does Medicare cover
physical therapy after knee replacement in Texas?''} must manually search through
hundreds of documents, cross-reference national and state-level policies, and interpret
dense regulatory prose. This process is time-consuming, error-prone, and inaccessible
to non-specialists. The consequences of misinterpretation are significant: incorrect
billing leads to claim denials, and uninformed patients may forgo covered treatments.
 
Large Language Models (LLMs) offer a promising path toward automating policy question
answering. However, general-purpose LLMs suffer from a critical limitation in this
context: \textit{hallucination} --- the tendency to generate plausible-sounding but
factually incorrect information. In the healthcare domain, a hallucinated coverage
determination can have serious real-world consequences.
 
Retrieval-Augmented Generation (RAG) addresses this by grounding the model's responses
in retrieved documents from a curated knowledge base. Rather than relying on parametric
knowledge from pre-training, a RAG system retrieves relevant policy chunks at query time
and passes them to the LLM as context, instructing it to generate answers only from
the provided evidence. This dramatically reduces hallucination while enabling precise,
citable responses.
 
This project implements a complete RAG pipeline for Medicare policy Q\&A using real CMS
data, evaluates it against an ungrounded LLM baseline, and delivers an interactive demo
via a Streamlit interface.
 
% ─────────────────────────────────────────────────────────────────────────────
\section{Literature Review}
% ─────────────────────────────────────────────────────────────────────────────
 
\subsection{Retrieval-Augmented Generation}
 
Retrieval-Augmented Generation was formally introduced by Lewis et al.\ (2020) as a
general framework for knowledge-intensive NLP tasks \cite{lewis2020rag}. The key insight
is that combining a parametric memory (the LLM) with a non-parametric memory (a retrieval
index over documents) produces systems that are more factually accurate, more interpretable,
and more easily updated than fine-tuned models alone.
 
Subsequent work has refined each pipeline component. Karpukhin et al.\ (2020) demonstrated
that dense bi-encoder retrieval substantially outperforms BM25 sparse retrieval for
open-domain question answering \cite{karpukhin2020dpr}. Nogueira and Cho (2019) introduced
neural reranking via BERT-based cross-encoders, showing that a two-stage retrieve-then-rerank
pipeline significantly improves precision over single-stage retrieval \cite{nogueira2019reranker}.
This design --- fast approximate retrieval followed by accurate reranking --- has become the
standard for production RAG systems and directly informs our architecture.
 
\subsection{RAG in Healthcare}
 
The healthcare domain presents unique NLP challenges: terminology is highly specialized,
documents are long and structurally complex, and factual accuracy carries direct clinical
and legal implications. Xiong et al.\ (2024) demonstrated that RAG systems significantly
outperform standard LLMs on medical QA benchmarks, particularly on questions requiring
specific factual recall \cite{xiong2024}. Crucially, they found that retrieval quality ---
specifically whether the correct source document appeared in the top-$k$ candidates ---
was the dominant factor in system performance, motivating our emphasis on reranking.
 
Singhal et al.\ (2023) showed that domain-adapted LLMs with grounded generation achieve
expert-level clinical QA performance, but their approach required expensive fine-tuning
on proprietary data \cite{singhal2023medpalm}. Our work pursues a complementary strategy:
rather than adapting the model, we focus on retrieval quality and grounding via prompting,
making the system more accessible and easily updated as policies change.
 
Narayanan et al.\ (2023) applied RAG to health insurance document Q\&A and found that
chunking strategy had a significant effect on retrieval quality, with section-aware chunking
outperforming naive fixed-length splitting \cite{narayanan2023insurance}. This finding
directly informed our chunking design decisions (see Section~\ref{sec:chunking}).
 
\subsection{Embedding Models and Rerankers}
 
Xiao et al.\ (2023) introduced the BGE (BAAI General Embedding) model family, achieving
state-of-the-art performance on the MTEB benchmark across retrieval, reranking, and
semantic similarity tasks \cite{xiao2023bge}. BGE models are open-source, freely
available on HuggingFace, and optimized for asymmetric retrieval where queries and
documents differ in length --- exactly our setting. We adopt \texttt{BAAI/bge-base-en-v1.5}
for dense retrieval and \texttt{BAAI/bge-reranker-v2-m3} for cross-encoder reranking.
 
\subsection{Hallucination and Faithfulness}
 
Maynez et al.\ (2020) provide a taxonomy of hallucination distinguishing \textit{intrinsic}
errors (contradicting the source) from \textit{extrinsic} errors (introducing absent
information) \cite{maynez2020hallucination}. Min et al.\ (2023) showed that faithfulness
--- the degree to which a generated answer is supported by retrieved evidence --- can be
reliably measured using an LLM-as-judge methodology \cite{min2023faitheval}. We adopt
both human annotation and LLM-judge evaluation in our experiments.
 
% ─────────────────────────────────────────────────────────────────────────────
\section{Project Scope and Objectives}
% ─────────────────────────────────────────────────────────────────────────────
 
\subsection{Problem Statement}
 
We address the following problem: \textit{given a natural language question about Medicare
coverage, generate a factually accurate, grounded, and citable answer derived from official
CMS policy documents.}
 
Representative questions the system is designed to answer include:
 
\begin{itemize}[noitemsep, topsep=4pt]
    \item \textit{``Does Medicare cover acupuncture for chronic lower back pain?''}
    \item \textit{``Is home oxygen therapy covered for patients in Texas?''}
    \item \textit{``What are the criteria for cardiac rehabilitation coverage in Florida?''}
    \item \textit{``Does Medicare cover continuous glucose monitoring devices?''}
    \item \textit{``What wound care treatments are covered under Novitas?''}
\end{itemize}
 
\subsection{Objectives}
 
The specific objectives of this project are:
 
\begin{enumerate}[noitemsep, topsep=4pt]
    \item Build a curated knowledge base from real CMS NCD and LCD documents.
    \item Implement a two-stage retrieval pipeline: dense embedding retrieval followed
          by cross-encoder reranking.
    \item Integrate retrieved evidence with an open-source LLM to produce grounded,
          cited answers.
    \item Evaluate the system on retrieval quality, answer correctness, hallucination
          rate, and citation accuracy.
    \item Compare three configurations: Baseline LLM, RAG, and RAG with query rewriting.
    \item Deliver a functional interactive demo via Streamlit.
\end{enumerate}
 
\subsection{Scope Boundaries}
 
This project focuses exclusively on Medicare coverage policy. We do not address Medicaid,
private insurance, or clinical treatment guidelines. The knowledge base is restricted to
NCDs and LCDs sourced from CMS as of early 2026. The system is a decision-support tool
and does not constitute medical or legal advice.
 
% ─────────────────────────────────────────────────────────────────────────────
\section{Data Collection and Preprocessing}
% ─────────────────────────────────────────────────────────────────────────────
 
\subsection{Overview of Data Sources}
 
All data was sourced from the \textbf{Centers for Medicare \& Medicaid Services (CMS)},
the federal agency responsible for administering Medicare. We collected four distinct
document corpora, each serving a different role in answering Medicare coverage questions.
Two were obtained as structured bulk exports from the CMS Medicare Coverage Database:
 
\begin{center}
    \url{https://www.cms.gov/medicare-coverage-database/downloads/downloads.aspx}
\end{center}
 
Two additional sources --- the Medicare Benefit Policy Manual and the Medicare Claims
Processing Manual --- were obtained as PDF chapter downloads from the CMS Internet-Only
Manuals repository:
 
\begin{center}
    \url{https://www.cms.gov/Regulations-and-Guidance/Guidance/Manuals/Internet-Only-Manuals-IOMs}
\end{center}
 
Together, these four sources provide comprehensive coverage of Medicare policy at both
the determination level (what is covered) and the operational level (how coverage is
defined and administered). Each source is described in detail below.
 
\subsection{National Coverage Determinations (NCDs)}
 
\subsubsection{What NCDs Are}
 
A National Coverage Determination is a federal-level ruling by CMS on whether a
particular item, service, treatment, or technology is covered under Medicare --- and if
so, under what clinical conditions. NCDs are binding across all 50 states and all
Medicare Administrative Contractors. They are issued after a formal evidence review
process and represent the most authoritative statement of Medicare coverage for a given
topic.
 
Each NCD is identified by a section number in the \textit{Medicare National Coverage
Determinations Manual} (CMS Publication 100-03), for example:
 
\begin{itemize}[noitemsep, topsep=4pt]
    \item \textbf{NCD 30.1} --- Acupuncture for Chronic Low Back Pain (added 2020)
    \item \textbf{NCD 100.3} --- 24-Hour Ambulatory Esophageal pH Monitoring
    \item \textbf{NCD 160.18} --- Continuous Glucose Monitors
    \item \textbf{NCD 20.4} --- Implantable Cardiac Pacemakers
    \item \textbf{NCD 80.12} --- Positron Emission Tomography (PET) Scans
\end{itemize}
 
A typical NCD specifies the clinical indications under which coverage applies, any
limitations or exclusions, and references the clinical evidence supporting the
determination. For example, NCD 30.1 (Acupuncture) covers up to 12 visits within
90 days for chronic low back pain lasting 12 or more weeks, with an additional 8
sessions if the patient demonstrates measurable improvement.
 
\subsubsection{How We Downloaded NCD Data}
 
The full NCD dataset was obtained as a structured bulk export from the CMS Downloads
page. The export consisted of four CSV files:
 
\begin{itemize}[noitemsep, topsep=4pt]
    \item \texttt{ncd\_trkg.csv} --- the primary table containing all NCD records with
          full HTML-formatted policy text
    \item \texttt{ncd\_pblctn\_ref.csv} --- a reference table mapping publication codes
          to manual section names
    \item \texttt{ncd\_bnft\_ctgry\_ref.csv} --- benefit category reference (e.g.,
          ``Home Health'', ``Durable Medical Equipment'')
    \item \texttt{ncd\_trkg\_bnft\_xref.csv} --- a cross-reference linking each NCD
          to its benefit categories
\end{itemize}
 
We parsed \texttt{ncd\_trkg.csv} as the primary source, filtering to active records
(\texttt{status = 'A'}) and removing duplicates by retaining the most recent version
of each NCD. After filtering, we extracted 343 unique active NCDs. The policy text
was HTML-formatted; we used Python's \texttt{BeautifulSoup} library to strip all HTML
tags and extract clean plain text. Each NCD was saved as a structured plain-text file:
 
\begin{verbatim}
TITLE: Acupuncture for Chronic Low Back Pain
SECTION: 30.1
EFFECTIVE DATE: 2020-01-21
COVERAGE POLICY:
Medicare covers acupuncture for chronic low back pain (cLBP)
lasting 12 weeks or longer; not associated with surgery...
\end{verbatim}
 
\subsection{Local Coverage Determinations (LCDs)}
 
\subsubsection{What LCDs Are}
 
While NCDs establish national-level policy, \textit{Local Coverage Determinations} are
issued by Medicare Administrative Contractors (MACs) --- private companies contracted
by CMS to administer Medicare claims in specific geographic regions. Each MAC issues
LCDs applicable to the states it serves, meaning coverage for a given procedure may
differ by location even when no national NCD exists.
 
The United States is divided into several MAC jurisdictions. Of particular relevance to
our corpus is \textbf{Novitas Solutions}, which serves Jurisdiction H (Texas, Arkansas,
Colorado, Louisiana, Mississippi, New Mexico, and Oklahoma) and Jurisdiction L
(Delaware, Maryland, New Jersey, Pennsylvania, and Washington DC). A user asking
\textit{``Does Medicare cover wound care in Texas?''} may receive a more precise answer
from a Novitas LCD than from any federal NCD. Our system preserves and leverages this
state metadata throughout the pipeline.
 
\subsubsection{Examples of LCDs in Our Corpus}
 
\begin{table}[h]
\centering
\caption{Sample LCD Documents in the Corpus}
\label{tab:lcdsample}
\begin{tabular}{@{}llll@{}}
\toprule
\textbf{LCD ID} & \textbf{Title} & \textbf{Contractor} & \textbf{States (sample)} \\
\midrule
L35125 & Wound Care                  & Novitas  & TX, AR, CO, LA \\
L33797 & Oxygen and Oxygen Equipment & CGS      & KY, OH \\
L37792 & 4Kscore Test Algorithm      & National & All states \\
L35062 & Physical Therapy            & Novitas  & TX, AR, CO \\
L37387 & Chiropractic Services       & WPS      & WI, MI, MN \\
L35050 & Outpatient Sleep Studies    & Palmetto & SC, NC, VA \\
\bottomrule
\end{tabular}
\end{table}
 
\subsubsection{How We Downloaded LCD Data}
 
The LCD bulk dataset was downloaded from the CMS Downloads page as a single Excel file
(\texttt{lcd\_half.xlsx}). The file presented a significant formatting challenge: the
HTML-encoded policy text had overflowed across hundreds of Excel columns during export,
producing a table with over 1,300 columns that standard tools could not parse correctly.
 
We resolved this programmatically using Python's \texttt{openpyxl} library, which
handles malformed column structures gracefully. After filtering to valid LCD records
(numeric \texttt{lcd\_id}, status \texttt{'A'}, and title length under 300 characters)
and deduplicating by LCD ID to retain the most recent version of each policy, we
obtained \textbf{819 active LCD records}.
 
A critical preprocessing step was \textbf{state attribution}: mapping each LCD to the
states it governs based on CMS's official MAC jurisdiction assignments. Each file header
was enriched with contractor and state metadata:
 
\begin{verbatim}
TITLE: Wound Care
LCD_ID: 35125
CONTRACTOR: Novitas
STATES: TX, AR, CO, LA, MS, NM, OK, DE, DC, MD, NJ, PA
EFFECTIVE_DATE: 2018-04-19
 
COVERAGE POLICY:
Medicare covers wound care services when the patient presents
with a chronic wound that has not responded to standard
treatment after 30 days...
\end{verbatim}
 
This state metadata is preserved in every downstream chunk, enabling state-aware
retrieval --- so a Texas-specific query automatically prioritizes chunks tagged
with \texttt{TX}.
 
\subsection{Medicare Benefit Policy Manual (Publication 100-02)}
 
\subsubsection{What the Benefit Policy Manual Is}
 
The \textit{Medicare Benefit Policy Manual} is CMS Publication 100-02 and serves as
the structural foundation for understanding what Medicare covers and under what
conditions. While NCDs and LCDs answer ``is procedure X covered?'', the Benefit Policy
Manual answers the more fundamental question: ``what does Medicare coverage actually
mean?'' It defines core concepts such as medical necessity, skilled nursing care,
homebound status, and the conditions under which different benefit categories apply.
 
The manual is organized into 17 chapters, each covering a distinct benefit category:
 
\begin{itemize}[noitemsep, topsep=4pt]
    \item \textbf{Chapter 1} --- Inpatient Hospital Services
    \item \textbf{Chapter 6} --- Hospital Services Covered Under Part B
    \item \textbf{Chapter 7} --- Home Health Services
    \item \textbf{Chapter 8} --- Skilled Nursing Facility (SNF) Coverage
    \item \textbf{Chapter 9} --- Hospice Services
    \item \textbf{Chapter 10} --- Ambulance Services
    \item \textbf{Chapter 11} --- End Stage Renal Disease (ESRD)
    \item \textbf{Chapter 15} --- Covered Medical and Other Health Services
    \item \textbf{Chapter 16} --- General Exclusions from Coverage
    \item \textit{... and 8 additional chapters}
\end{itemize}
 
Chapter 16 (General Exclusions from Coverage) is particularly valuable for our system,
as it enumerates categories of services that Medicare explicitly does not cover ---
enabling the system to correctly answer negative questions such as \textit{``Does
Medicare cover routine dental care?''} (it does not).
 
\subsubsection{How We Downloaded the Benefit Policy Manual}
 
All 17 chapters were downloaded as individual PDF files from the CMS Internet-Only
Manuals repository. Text was extracted from each PDF using the \texttt{PyMuPDF}
(\texttt{fitz}) library, which preserves paragraph structure and handles the
multi-column layouts common in CMS documents. Each chapter was tagged with its
manual name, chapter title, and a scope of \texttt{STATES: ALL} since benefit policy
applies nationally.
 
\subsection{Medicare Claims Processing Manual (Publication 100-04)}
 
\subsubsection{What the Claims Processing Manual Is}
 
The \textit{Medicare Claims Processing Manual} is CMS Publication 100-04 and governs
how Medicare claims are submitted, processed, and adjudicated. While primarily
administrative in nature, it contains substantial coverage-relevant content: it
specifies billing codes, documentation requirements, and the exact conditions under
which specific services are reimbursable. For a RAG system answering questions from
healthcare providers, this manual is important because providers frequently need to
know not just whether something is covered, but what documentation is required to
support a claim.
 
The manual spans 39 chapters covering the full range of Medicare claim types:
 
\begin{itemize}[noitemsep, topsep=4pt]
    \item \textbf{Chapter 1} --- General Billing Requirements
    \item \textbf{Chapter 10} --- Home Health Agency Billing
    \item \textbf{Chapter 11} --- Processing Hospice Claims
    \item \textbf{Chapter 12} --- Physicians and Nonphysician Practitioners
    \item \textbf{Chapter 16} --- Laboratory Services
    \item \textbf{Chapter 18} --- Preventive and Screening Services
    \item \textbf{Chapter 20} --- Durable Medical Equipment (DMEPOS)
    \item \textit{... and 32 additional chapters}
\end{itemize}
 
\subsubsection{How We Downloaded the Claims Processing Manual}
 
All 39 chapters were downloaded as individual PDF files from the same CMS Internet-Only
Manuals repository and processed identically to the Benefit Policy Manual --- PDF text
extraction via \texttt{PyMuPDF}, chunked at 400 words with 50-word overlap, and tagged
with \texttt{MANUAL: Medicare Claims Processing Manual}, chapter title, and
\texttt{STATES: ALL}.
 
\subsection{Final Corpus Summary}
 
\begin{table}[h]
\centering
\caption{Complete Data Collection Summary}
\label{tab:datasources}
\begin{tabular}{@{}llll@{}}
\toprule
\textbf{Document Type} & \textbf{Format} & \textbf{Source} & \textbf{Count} \\
\midrule
National Coverage Determinations  & CSV bulk export  & CMS Coverage Database & 343 docs \\
Local Coverage Determinations     & Excel bulk export & CMS Coverage Database & 819 docs \\
Medicare Benefit Policy Manual    & PDF (per chapter) & CMS IOMs & 17 chapters \\
Medicare Claims Processing Manual & PDF (per chapter) & CMS IOMs & 39 chapters \\
\midrule
\textbf{Total} & & & \textbf{1,218 documents} \\
\bottomrule
\end{tabular}
\end{table}
 
\noindent The four sources are complementary by design. NCDs and LCDs answer specific
coverage questions at the policy level. The Benefit Policy Manual provides the structural
definitions and benefit rules that underpin those policies. The Claims Processing Manual
adds the provider-facing operational layer: documentation requirements, billing codes,
and submission criteria. Together they form a knowledge base that can answer the full
range of Medicare coverage questions a provider or patient might ask.
 
% ─────────────────────────────────────────────────────────────────────────────
\section{Step 1: Document Chunking and Corpus Construction}
\label{sec:chunking}
% ─────────────────────────────────────────────────────────────────────────────
 
\subsection{Why Chunking Is Necessary}
 
A fundamental challenge in building a retrieval system over long documents is that
embedding models encode a fixed-length vector for each input. A full NCD or LCD
document may contain 1,000--3,000 words: if the entire document is encoded as a
single vector, semantically distinct sections average together and retrieval precision
degrades. Similarly, passing complete documents to an LLM at generation time is
inefficient and may exceed context window limits.
 
The standard solution is to split documents into smaller, overlapping \textit{chunks}
before indexing. Each chunk is independently embedded and stored. At query time, only
the most relevant chunks are retrieved and passed to the LLM --- not entire documents.
 
\subsection{Chunking Strategy and Parameters}
 
We adopted a \textbf{fixed-size sliding window} approach with the following parameters:
 
\begin{itemize}[noitemsep, topsep=4pt]
    \item \textbf{Chunk size}: 400 words
    \item \textbf{Overlap}: 50 words between consecutive chunks
\end{itemize}
 
The 400-word size balances retrieval granularity and semantic coherence. Chunks shorter
than 200 words often lack sufficient context for the embedding model to capture policy
meaning. Chunks longer than 600 words begin to mix multiple distinct coverage criteria,
reducing retrieval precision. The 50-word overlap ensures that sentences near chunk
boundaries appear in at least two chunks, preventing information loss at split points.
 
\subsection{Metadata Tagging Design}
 
A critical design decision was to embed document metadata \textit{directly into every
chunk's text}, rather than storing it only as a separate database field. This ensures
the LLM always has access to citation information regardless of how the chunk is
retrieved. Every chunk is prefixed with a structured metadata header:
 
For NCD chunks:
\begin{verbatim}
[TYPE: NCD | NCD_ID: 30.1 | TITLE: Acupuncture for Chronic Low Back Pain]
\end{verbatim}
 
For LCD chunks:
\begin{verbatim}
[TYPE: LCD | LCD_ID: 35125 | TITLE: Wound Care |
 STATES: TX, AR, CO, LA | CONTRACTOR: Novitas]
\end{verbatim}
 
This allows the LLM to generate citations such as \textit{``According to LCD L35125
(Wound Care, Novitas Solutions)...''} directly from the chunk content, without any
additional lookup.
 
\subsection{Implementation}
 
The chunking pipeline was implemented in Python and executed in Google Colab Pro.
The core chunking function is:
 
\begin{lstlisting}[language=Python, caption=Core sliding-window chunking function]
def chunk_text(text, chunk_size=400, overlap=50):
    words = text.split()
    chunks = []
    i = 0
    while i < len(words):
        chunk = ' '.join(words[i:i + chunk_size])
        chunks.append(chunk)
        i += chunk_size - overlap
        if i + overlap >= len(words):
            break
    return chunks
\end{lstlisting}
 
Metadata is parsed from the structured file headers and prepended to each chunk.
The LCD parser, which also extracts state information, is shown below:
 
\begin{lstlisting}[language=Python, caption=LCD file parser with state metadata extraction]
def parse_lcd_file(filepath, filename):
    with open(filepath, encoding='utf-8') as f:
        content = f.read()
 
    lines = content.split('\n')
    title, lcd_id, contractor = '', 'UNKNOWN', 'UNKNOWN'
    states = ['ALL']
 
    for line in lines[:15]:
        if line.startswith('TITLE:'):
            title = line.replace('TITLE:', '').strip()
        elif line.startswith('LCD_ID:'):
            lcd_id = line.replace('LCD_ID:', '').strip()
        elif line.startswith('CONTRACTOR:'):
            contractor = line.replace('CONTRACTOR:', '').strip()
        elif line.startswith('STATES:'):
            raw = line.replace('STATES:', '').strip()
            states = [s.strip() for s in raw.split(',')]
 
    text_chunks = chunk_text(content)
    results = []
 
    for i, chunk in enumerate(text_chunks):
        header = (
            f"[TYPE: LCD | LCD_ID: {lcd_id} | TITLE: {title}"
            f" | STATES: {', '.join(states)}"
            f" | CONTRACTOR: {contractor}]"
        )
        tagged_chunk = f"{header}\n\n{chunk}"
        results.append({
            'text':       tagged_chunk,
            'raw_text':   chunk,
            'source_id':  lcd_id,
            'title':      title,
            'type':       'LCD',
            'states':     states,
            'contractor': contractor,
            'filename':   filename,
            'chunk_idx':  i
        })
    return results
\end{lstlisting}
 
All 1,218 processed documents were merged into a single \texttt{all\_chunks.json} file,
where each entry follows this schema:
 
\begin{lstlisting}[language=Python, caption=JSON schema for each chunk in all\_chunks.json]
{
    "text":       "[TYPE: LCD | LCD_ID: 35125 | ...]\n\nMedicare covers...",
    "raw_text":   "Medicare covers wound care when the patient...",
    "source_id":  "35125",
    "title":      "Wound Care",
    "type":       "LCD",
    "states":     ["TX", "AR", "CO", "LA"],
    "contractor": "Novitas",
    "filename":   "Wound_Care_35125.txt",
    "chunk_idx":  0
}
\end{lstlisting}
 
\subsection{Corpus Statistics}
 
\begin{table}[h]
\centering
\caption{Corpus Statistics After Chunking}
\label{tab:corpus}
\begin{tabular}{@{}lrrr@{}}
\toprule
\textbf{Document Type} & \textbf{Files} & \textbf{Chunks} & \textbf{Avg.\ Chunks / Doc} \\
\midrule
National Coverage Determinations  & 343  & 575   & 1.68 \\
Local Coverage Determinations     & 819  & 2,802 & 3.42 \\
Medicare Benefit Policy Manual    & 17   & 1,039 & 61.1  \\
Medicare Claims Processing Manual & 39   & 4,147 & 106.3 \\
\midrule
\textbf{Total}                    & \textbf{1,218} & \textbf{8,563} & \textbf{7.03} \\
\bottomrule
\end{tabular}
\end{table}
 
The substantially higher chunk counts for the two CMS manuals reflect their greater
length and scope: each PDF chapter spans dozens of pages of detailed policy, billing
codes, and administrative guidance, compared to the more concise NCD and LCD documents.
The complete corpus of 8,563 chunks in \texttt{all\_chunks.json} serves as input to
all downstream pipeline components.
 

% ─────────────────────────────────────────────────────────────────────────────
\section{Step 2: Embeddings and Vector Index}
% ─────────────────────────────────────────────────────────────────────────────

\subsection{Embedding Model Selection}

To convert text chunks into dense vector representations suitable for semantic search,
we selected \texttt{BAAI/bge-base-en-v1.5} from the BGE (BAAI General Embedding) model
family. This model was chosen for several reasons:

\begin{itemize}[noitemsep, topsep=4pt]
    \item \textbf{Retrieval benchmark performance}: BGE models achieve state-of-the-art
          scores on the Massive Text Embedding Benchmark (MTEB) across retrieval,
          reranking, and semantic similarity tasks, outperforming other base-sized models
          at the time of selection.
    \item \textbf{Asymmetric retrieval optimization}: BGE is specifically designed for
          settings where queries are short (user questions) and documents are longer
          (policy chunks). This asymmetry is a poor fit for models trained primarily 
          on symmetric sentence similarity.
    \item \textbf{768-dimensional output}: The base variant produces 768-dimensional
          embeddings, providing a strong balance between representation quality and
          computational efficiency. Larger variants (e.g., \texttt{bge-large}) offer
          marginal gains at significantly higher cost.
    \item \textbf{Open-source availability}: The model is freely available on HuggingFace
          under the MIT license, with no API costs which is important for a reproducible
          academic project.
\end{itemize}

\subsection{Encoding the Corpus}

All 8,563 chunks from \texttt{all\_chunks.json} were encoded into dense vectors using
the \texttt{sentence-transformers} library. We embedded the \textbf{metadata-tagged
\texttt{text} field} (not the \texttt{raw\_text} field) from each chunk. This is an
intentional design decision: the metadata header (e.g.,
\texttt{[TYPE: LCD | LCD\_ID: 35125 | TITLE: Wound Care | STATES: TX, AR, CO, LA |
CONTRACTOR: Novitas]}) is semantically informative. By embedding it alongside the
policy content, a query about ``wound care Texas'' can match on both the policy text
and the metadata tags, improving retrieval quality.

A critical encoding parameter is \textbf{L2 normalization}: each embedding vector is
normalized to unit length during encoding (\texttt{normalize\_embeddings=True}). This
ensures that the inner product between any two vectors equals their cosine similarity,
which simplifies the FAISS index configuration (see below). After encoding, we verified
that all 8,563 vectors had an L2 norm of exactly 1.0.

The encoding was performed on a GPU-enabled Google Colab instance with a batch size
of 64. The core encoding logic is:

\begin{lstlisting}[language=Python, caption=Encoding all chunks into normalized dense vectors]
from sentence_transformers import SentenceTransformer
import numpy as np

model = SentenceTransformer('BAAI/bge-base-en-v1.5', device='cuda')

texts = [chunk['text'] for chunk in all_chunks]

embeddings = model.encode(
    texts,
    batch_size=64,
    show_progress_bar=True,
    normalize_embeddings=True,
    convert_to_numpy=True
)
# embeddings.shape: (8563, 768), dtype: float32
\end{lstlisting}

\subsection{FAISS Index Construction}

We use \textbf{FAISS} (Facebook AI Similarity Search) to index the embedding vectors
for fast nearest-neighbor retrieval. FAISS is a library developed by Meta AI for
efficient similarity search over dense vectors, supporting both exact and approximate
search algorithms.

We selected \texttt{IndexFlatIP} (flat inner product index) for the following reasons:

\begin{enumerate}[noitemsep, topsep=4pt]
    \item \textbf{Exact search}: With approximately 8,500 vectors, brute-force search
          completes in 1--5ms per query which is fast enough for interactive use. Approximate
          methods (IVF, HNSW) add complexity without meaningful speedup at this scale.
    \item \textbf{Cosine similarity via inner product}: Since all embeddings are
          L2-normalized, the inner product between any two vectors equals their cosine
          similarity. This avoids a separate normalization step at query time.
    \item \textbf{No training required}: Unlike approximate indexes such as
          \texttt{IndexIVFFlat}, \texttt{IndexFlatIP} requires no training or clustering
          step where vectors are simply added to the index.
    \item \textbf{Deterministic results}: Every query returns the mathematically exact
          top-$k$ results, ensuring full reproducibility.
\end{enumerate}

The index construction code is straightforward:

\begin{lstlisting}[language=Python, caption=Building and saving the FAISS index]
import faiss

dimension = embeddings.shape[1]  # 768
index = faiss.IndexFlatIP(dimension)
index.add(embeddings.astype('float32'))

# Save to Google Drive
faiss.write_index(index, 'faiss_index/medicare.index')
\end{lstlisting}

\subsection{Index Statistics}

\begin{table}[h]
\centering
\caption{Embedding and Index Statistics}
\label{tab:indexstats}
\begin{tabular}{@{}lr@{}}
\toprule
\textbf{Metric} & \textbf{Value} \\
\midrule
Total chunks indexed        & 8,563 \\
Embedding model             & \texttt{BAAI/bge-base-en-v1.5} \\
Embedding dimension         & 768 \\
Normalization               & L2-normalized \\
Similarity metric           & Cosine (via inner product) \\
FAISS index type            & \texttt{IndexFlatIP} (exact) \\
Index file size             & 25.1 MB \\
Embeddings file size        & 25.1 MB \\
Metadata file size          & 4.0 MB \\
\bottomrule
\end{tabular}
\end{table}

Four artifacts were saved to Google Drive for use by downstream pipeline components:
the FAISS index file (\texttt{medicare.index}), a document-ID mapping
(\texttt{docids.txt}), the raw NumPy embedding array (\texttt{embeddings.npy}), and
a chunk metadata JSON file (\texttt{chunk\_metadata.json}) containing the type, title,
states, contractor, and filename for each indexed chunk.

\subsection{Retrieval Sanity Test}

Before handing off to the reranking pipeline, we verified that the FAISS index
retrieves semantically relevant chunks by running three sample queries spanning
different document types. For each query, the top 3 results (before any reranking)
are shown below.

\textbf{Query 1} (NCD --- national coverage): \textit{``Does Medicare cover acupuncture
for chronic lower back pain?''}

\begin{table}[h]
\centering
\small
\begin{tabular}{@{}clllc@{}}
\toprule
\textbf{Rank} & \textbf{Type} & \textbf{ID} & \textbf{Title} & \textbf{Score} \\
\midrule
1 & NCD & 30.3.3 & Acupuncture for Chronic Lower Back Pain (cLBP) & 0.8095 \\
2 & NCD & 30.3   & Acupuncture                                    & 0.7742 \\
3 & NCD & 30.3.1 & Acupuncture for Fibromyalgia                   & 0.7358 \\
\bottomrule
\end{tabular}
\end{table}

The top result is the exact NCD governing acupuncture for chronic lower back pain,
confirming that the embedding model captures the semantic match between the query
and the policy document.

\textbf{Query 2} (LCD --- state-specific): \textit{``Is home oxygen therapy covered
for patients in Texas?''}

\begin{table}[h]
\centering
\small
\begin{tabular}{@{}clllc@{}}
\toprule
\textbf{Rank} & \textbf{Type} & \textbf{ID} & \textbf{Title} & \textbf{Score} \\
\midrule
1 & NCD & 240.2   & Home Use of Oxygen                    & 0.8005 \\
2 & NCD & 240.2   & Home Use of Oxygen (chunk 2)          & 0.7805 \\
3 & NCD & 240.2.1 & Home Use of Oxygen in Clinical Trials & 0.7583 \\
\bottomrule
\end{tabular}
\end{table}

The retrieval correctly identifies the NCD governing home oxygen use. Notably, the
LCD for ``Oxygen and Oxygen Equipment'' (LCD 33797, which covers Texas) appears at
rank 4 with a score of 0.7432. This illustrates a limitation of dense retrieval alone:
the NCD is semantically closer to the query, but for a Texas-specific question the LCD
may be more operationally relevant. State filtering and reranking (Step 3) address this.

\textbf{Query 3} (Benefit Policy): \textit{``What are the requirements for skilled
nursing facility coverage?''}

\begin{table}[h]
\centering
\small
\begin{tabular}{@{}clllc@{}}
\toprule
\textbf{Rank} & \textbf{Type} & \textbf{ID} & \textbf{Title} & \textbf{Score} \\
\midrule
1 & Benefit Policy & Ch.\ 7  & Home Health Services           & 0.7718 \\
2 & Benefit Policy & Ch.\ 8  & Coverage of Extended Care (SNF) & 0.7687 \\
3 & Benefit Policy & Ch.\ 8  & Coverage of Extended Care (SNF) & 0.7645 \\
\bottomrule
\end{tabular}
\end{table}

Chapter 8 (SNF Coverage) appears at ranks 2 and 3, confirming that the embeddings
successfully capture the relationship between ``skilled nursing facility'' in the query
and ``extended care (SNF)'' in the manual chapter titles. Chapter 7 (Home Health) at
rank 1 reflects the semantic overlap between home health and SNF eligibility criteria
--- both involve skilled care requirements.

These results demonstrate that the FAISS index provides strong semantic retrieval
across all four document types in the corpus. The next pipeline stage (Step 3) adds
state-aware filtering and cross-encoder reranking to further improve precision.


\newpage
% ─────────────────────────────────────────────────────────────────────────────
\section{Step 3: Retrieval Pipeline and Reranking}
% ─────────────────────────────────────────────────────────────────────────────

In order to retrieve the most relevant Medicare policy information for a user query, MediQuery implements a two-stage retrieval pipeline commonly used in modern Retrieval-Augmented Generation (RAG) systems. This approach combines fast semantic retrieval using dense embeddings with a more precise neural reranking step.

\subsection{Two-Stage Retrieve-Then-Rerank Architecture}

The retrieve-then-rerank architecture is designed to balance efficiency and accuracy. Searching across thousands of document chunks using a highly accurate neural model would be computationally expensive. Instead, the system first performs fast approximate retrieval to identify a small candidate set of potentially relevant documents. A more computationally intensive model is then applied to refine the ranking of those candidates.

The architecture used in MediQuery follows the pipeline shown below:

\begin{center}
Query $\rightarrow$ Embedding Model $\rightarrow$ FAISS Retrieval $\rightarrow$ Candidate Chunks $\rightarrow$ Cross-Encoder Reranker $\rightarrow$ Top Evidence Chunks $\rightarrow$ LLM
\end{center}

This architecture ensures that the system remains responsive while still providing highly relevant evidence for answer generation.

\subsection{Stage 1: Dense Retrieval}

The first stage performs semantic retrieval using dense vector embeddings. When a user submits a question, the system encodes the query using the same embedding model used to encode the document corpus, \texttt{BAAI/bge-base-en-v1.5}. This ensures that both queries and documents are represented in the same semantic vector space.

Once the query embedding is generated, the system performs nearest-neighbor search in the FAISS index containing all document embeddings. The index used in this project is \texttt{IndexFlatIP}, which performs similarity search using inner product operations. Because all embeddings are normalized, inner product corresponds to cosine similarity.

The FAISS search returns the top 20 document chunks that are most similar to the query embedding. These chunks form the initial candidate set of potentially relevant evidence.

\subsection{State-Aware Filtering}

Medicare policies often vary by geographic region because Local Coverage Determinations (LCDs) are issued by Medicare Administrative Contractors that cover specific states. During preprocessing, each document chunk was tagged with metadata including the states covered by that policy.

After retrieving the candidate chunks from FAISS, the system applies a state-aware filtering step. If a user query mentions a specific state (for example Texas), the system prioritizes chunks that contain the corresponding state code in their metadata. For example, a query related to Texas will prioritize chunks tagged with \texttt{TX}. This filtering step ensures that the system retrieves policies that are relevant to the user's geographic context.

\subsection{Stage 2: Cross-Encoder Reranking}

While dense retrieval is efficient, it can sometimes return documents that are semantically related but not directly relevant to the user's question. To address this limitation, MediQuery applies a neural reranking model to the candidate set.

The reranking model used in this project is \texttt{BAAI/bge-reranker-v2-m3}. This model evaluates the relevance of each query-document pair by jointly processing the query and the document text.

For each retrieved chunk, the system constructs a pair consisting of the query and the chunk text. The reranker assigns a relevance score to each pair. The candidate chunks are then sorted according to their reranker scores, and the top five chunks are selected as the final evidence set.

\subsection{Bi-Encoder vs Cross-Encoder Models}

Dense retrieval models such as BGE embeddings use a bi-encoder architecture. In a bi-encoder, the query and document are encoded independently into vectors, and similarity is computed using vector operations. This approach is very fast and allows pre-computation of document embeddings, making it suitable for large-scale retrieval.

However, because the query and document are encoded separately, the model cannot fully capture detailed interactions between the two pieces of text. As a result, bi-encoder retrieval can sometimes produce approximate rankings.

Cross-encoders address this limitation by processing the query and document together as a single input sequence. This allows the model to capture richer contextual relationships between the query and the document text. The trade-off is that cross-encoders are significantly slower because they must evaluate each query-document pair individually.

By combining both approaches, the MediQuery system benefits from the speed of bi-encoder retrieval and the accuracy of cross-encoder reranking.

\subsection{Retrieve and Rerank Implementation}

The retrieval pipeline is implemented using a Python function that retrieves candidate chunks from the FAISS index and then reranks them using the cross-encoder model.

\begin{lstlisting}[language=Python]
def rerank_results(query, results, top_n=5):

    pairs = [(query, r["text"]) for r in results]

    scores = reranker.predict(pairs)

    for i in range(len(results)):
        results[i]["rerank_score"] = float(scores[i])

    results.sort(
        key=lambda x: x["rerank_score"],
        reverse=True
    )

    seen = set()
    final = []

    for r in results:
        if r["source_id"] not in seen:
            final.append(r)
            seen.add(r["source_id"])

        if len(final) == top_n:
            break

    return final
\end{lstlisting}

This function computes relevance scores for each retrieved chunk and returns the top ranked chunks after removing duplicate source documents.

\subsection{Before and After Reranking Comparison}

\begin{table}[h]
\centering
\begin{tabular}{p{4cm} p{5cm} p{5cm}}
\toprule
\textbf{Query} & \textbf{Top Result Before Reranking} & \textbf{Top Result After Reranking} \\
\midrule
Does Medicare cover acupuncture for chronic lower back pain? &
General pain management guideline related to acupuncture &
Medicare policy section covering acupuncture for chronic lower back pain \\
\addlinespace
Does Medicare cover wound care in Texas? &
National wound care overview document &
Texas Local Coverage Determination (LCD) policy for wound care services \\
\bottomrule
\end{tabular}
\caption{Example comparison of retrieval results before and after reranking}
\end{table}

To evaluate the impact of reranking, we compared the top retrieved documents before and after applying the cross-encoder model.  

The results show that reranking improves retrieval precision by promoting policy chunks that more directly answer the user's question.

\subsection{Fine-Tuning the Bi-Encoder}

While the pre-trained BGE embedding model provides strong general-purpose retrieval, its
representations are not optimized for the specialized vocabulary and structure of Medicare
policy documents. To close this domain gap, we fine-tuned both the bi-encoder and the
cross-encoder reranker on Medicare-specific training data derived from our corpus.

\subsubsection{Training Data Construction}

We constructed a training set from the 56,288 synthetic query-gold chunk pairs described
in Section~\ref{sec:chunking}. Each pair links a natural-language question to its
source chunk in \texttt{all\_chunks.json}. The positive passage for each training example
is the \textbf{full indexed chunk text} (including the metadata header), not the short
excerpt from the evaluation CSV. This ensures that the fine-tuned model learns
representations consistent with the actual documents it will search at inference time.

\subsubsection{Hard Negative Mining}

Effective contrastive learning requires not only positive pairs but also challenging
negative examples that the model currently confuses with true matches. For each of the
56,288 queries, we encoded the query with the pre-trained \texttt{bge-base-en-v1.5}
model and retrieved the top 300 candidate chunks from the pre-trained FAISS index. The
hard negative was selected as the \textbf{highest-ranked chunk whose source document ID
differs} from the gold chunk's source ID. This guarantees that each negative is
semantically close to the query but originates from a different policy
document. All 56,288 pairs received valid hard negatives.

\subsubsection{Train / Validation / Test Split}

The 56,288 triplets were split in two stages, both stratified by query type to maintain
label balance:

\begin{enumerate}[noitemsep, topsep=4pt]
    \item 10\% held out as a \textbf{test set} (5,629 examples) for final evaluation.
    \item The remaining 90\% split 95/5 into \textbf{train} (48,126) and \textbf{validation}
          (2,533).
\end{enumerate}

\subsubsection{Loss Function and Hyperparameters}

We fine-tuned \texttt{BAAI/bge-base-en-v1.5} using \textbf{MultipleNegativesRankingLoss}
(MNRL) with a scale factor of 50.0 (equivalent to a contrastive temperature of 0.02),
matching the temperature used in BGE's original pre-training. Each training batch
contained 32 unique anchors; the \texttt{NO\_DUPLICATES} batch sampler ensured maximum
in-batch negative diversity.

\begin{table}[h]
\centering
\caption{Bi-Encoder Fine-Tuning Hyperparameters}
\label{tab:biencoder-hp}
\begin{tabular}{@{}lr@{}}
\toprule
\textbf{Parameter} & \textbf{Value} \\
\midrule
Base model              & \texttt{BAAI/bge-base-en-v1.5} \\
Epochs                  & 2 \\
Batch size              & 32 \\
Learning rate           & $1 \times 10^{-5}$ \\
Warmup steps            & 300 (10\% of 3,006 total) \\
Temperature             & 0.02 (scale = 50) \\
Precision               & fp16 \\
Batch sampler           & \texttt{NO\_DUPLICATES} \\
Best-model metric       & \texttt{cosine\_ndcg@10} \\
\bottomrule
\end{tabular}
\end{table}

Validation was performed every 500 steps using an \texttt{InformationRetrievalEvaluator}
that searched 2,533 validation queries against the full 8,563-chunk corpus. The best
checkpoint was selected by cosine NDCG@10.

The core training setup is shown below:

\begin{lstlisting}[language=Python, caption=Bi-encoder fine-tuning with MNRL loss]
from sentence_transformers import SentenceTransformer
from sentence_transformers.losses import MultipleNegativesRankingLoss
from sentence_transformers.training_args import BatchSamplers

biencoder_model = SentenceTransformer('BAAI/bge-base-en-v1.5')
train_loss = MultipleNegativesRankingLoss(biencoder_model, scale=50.0)

args = SentenceTransformerTrainingArguments(
    output_dir=FT_BIENCODER_DIR,
    num_train_epochs=2,
    per_device_train_batch_size=32,
    learning_rate=1e-5,
    warmup_steps=300,
    fp16=True,
    batch_sampler=BatchSamplers.NO_DUPLICATES,
    load_best_model_at_end=True,
    metric_for_best_model="cosine_ndcg@10",
)
trainer = SentenceTransformerTrainer(
    model=biencoder_model, args=args,
    train_dataset=biencoder_train_dataset,
    loss=train_loss,
    evaluator=biencoder_evaluator,
)
trainer.train()
\end{lstlisting}

\subsubsection{Re-Encoding and Index Rebuild}

After fine-tuning, all 8,563 chunks were re-encoded with the fine-tuned bi-encoder
(\texttt{batch\_size=64}, \texttt{normalize\_embeddings=True}). A new FAISS
\texttt{IndexFlatIP} index was built from the fine-tuned embeddings and saved as
\texttt{medicare\_finetuned.index} (25.1\,MB).

\subsection{Fine-Tuning the Cross-Encoder Reranker}

\subsubsection{Hard Negative Re-Mining}

Because the fine-tuned bi-encoder produces a different retrieval ranking than the
pre-trained model, hard negatives for reranker training were \textbf{re-mined from the
fine-tuned FAISS index}. For each of the 48,126 training queries, we retrieved the top
100 candidates from the fine-tuned index and selected up to 7 hard negatives per query,
each from a distinct source document. This yields a training group size of 8 per query
(1~positive + 7~negatives), following the FlagEmbedding training protocol. After
filtering, the reranker training set contained 307,877 pairwise examples.

\subsubsection{Loss Function and Hyperparameters}

We fine-tuned \texttt{BAAI/bge-reranker-v2-m3} using \textbf{BinaryCrossEntropyLoss}
for pairwise relevance classification. Each example consisted of a query-chunk pair
labeled 1.0 (positive) or 0.0 (negative).

\begin{table}[h]
\centering
\caption{Reranker Fine-Tuning Hyperparameters}
\label{tab:reranker-hp}
\begin{tabular}{@{}lr@{}}
\toprule
\textbf{Parameter} & \textbf{Value} \\
\midrule
Base model              & \texttt{BAAI/bge-reranker-v2-m3} \\
Epochs                  & 2 \\
Batch size              & 32 \\
Learning rate           & $1 \times 10^{-5}$ \\
Warmup steps            & 10\% of total steps \\
Weight decay            & 0.01 \\
Precision               & fp16 \\
Gradient checkpointing  & Enabled \\
Max sequence length     & 2,048 tokens \\
Best-model metric       & \texttt{ndcg@5} \\
\bottomrule
\end{tabular}
\end{table}

Validation used a \texttt{CrossEncoderRerankingEvaluator} with 2,042 validation samples,
each containing 1~positive and its corresponding hard negatives. The best checkpoint was
selected by NDCG@5.

\subsection{Retrieval Improvement After Fine-Tuning}

To quantify the impact of domain-specific fine-tuning, we evaluated both the pre-trained
and fine-tuned pipelines on the held-out test set of 5,629 queries. FAISS retrieval
returned the top 20 candidates; the reranker selected the top 5 from the top 20 FAISS
results.

\begin{table}[h]
\centering
\caption{Retrieval Comparison: Pre-trained vs.\ Fine-tuned Pipeline}
\label{tab:ft-retrievalcompare}
\small
\begin{tabular}{@{}lccc@{}}
\toprule
\textbf{Metric} & \textbf{Pre-trained} & \textbf{Fine-tuned} & \textbf{Delta} \\
\midrule
Recall@5 (FAISS)     & 0.6460 & 0.7850 & +0.1390 \\
Recall@10 (FAISS)    & 0.7260 & 0.8420 & +0.1160 \\
Recall@20 (FAISS)    & 0.8040 & 0.9010 & +0.0970 \\
MRR (FAISS)          & 0.4992 & 0.6120 & +0.1128 \\
NDCG@5 (FAISS)       & 0.5224 & 0.6450 & +0.1226 \\
NDCG@10 (FAISS)      & 0.5488 & 0.6730 & +0.1242 \\
NDCG@20 (FAISS)      & 0.5686 & 0.6910 & +0.1224 \\
\midrule
Recall@5 (reranked)  & 0.6720 & 0.8100 & +0.1380 \\
MRR (reranked)       & 0.4751 & 0.6300 & +0.1549 \\
NDCG@5 (reranked)    & 0.5244 & 0.6650 & +0.1406 \\
\bottomrule
\end{tabular}
\end{table}

Fine-tuning produced consistent improvements across every metric. The largest
gain was in reranked MRR (+0.1549), indicating that fine-tuning not only increases whether
the correct chunk is retrieved but also how high it ranks among the final candidates shown
to the generator. Reranked Recall@5 reached 0.8100, meaning the correct evidence chunk
appears in the top-5 context for over 81\% of test queries.

\subsubsection{Impact of Reranking on the Fine-Tuned Pipeline}

Even after fine-tuning the bi-encoder, the cross-encoder reranker provides an additional
lift by re-scoring candidates with full query--document attention:

\begin{table}[h]
\centering
\caption{Fine-Tuned Pipeline: Before vs.\ After Reranking}
\label{tab:ft-rerank-delta}
\begin{tabular}{@{}lccc@{}}
\toprule
\textbf{Metric} & \textbf{FAISS Only} & \textbf{After Reranking} & \textbf{Delta} \\
\midrule
Recall@5  & 0.7850 & 0.8100 & +0.0250 \\
MRR       & 0.6120 & 0.6300 & +0.0180 \\
NDCG@5    & 0.6450 & 0.6650 & +0.0200 \\
\bottomrule
\end{tabular}
\end{table}

\subsubsection{Performance by Document Type}

Table~\ref{tab:ft-doctype} breaks down reranked Recall@5 by document corpus. Claims
Processing and LCD documents benefit most, likely because their larger chunk counts
provide more contrastive training signal.

\begin{table}[h]
\centering
\caption{Fine-Tuned Reranked Recall@5 by Document Type}
\label{tab:ft-doctype}
\begin{tabular}{@{}lrc@{}}
\toprule
\textbf{Document Type} & \textbf{Test Queries} & \textbf{Recall@5 (reranked)} \\
\midrule
Claims Processing Manual  & 2,535 & 0.8251 \\
Local Coverage Determinations & 1,944 & 0.8015 \\
Benefit Policy Manual     & 693   & 0.7958 \\
National Coverage Determinations & 457 & 0.7781 \\
\bottomrule
\end{tabular}
\end{table}

The final reranked chunks are then passed to the language model in Step 4, where they are used as evidence for generating grounded answers with citations.

\newpage

% ─────────────────────────────────────────────────────────────────────────────
\section{Step 4: LLM Generation and Citation}
% ─────────────────────────────────────────────────────────────────────────────

\subsection{Generator Choice and Role in the Pipeline}

The final stage of MediQuery transforms the top reranked evidence chunks into a concise,
patient-readable answer with explicit citations. For this component, we implemented a
Colab-friendly LLM synthesis module using \texttt{gemini-3.1-flash-} \texttt{lite-preview}. The
motivation for this choice was practical as well as technical: the model was easy to call
through a lightweight REST interface, responded quickly enough for interactive demos, and
produced stable outputs in a structured format without requiring local GPU inference.
Although the surrounding RAG pipeline is model-agnostic, Gemini provided an efficient way
to test grounded generation, citation formatting, and fallback behavior during integration.

The generator does not see the full corpus. Instead, it receives only the user's question
and the top five reranked chunks selected by the retrieval pipeline. This design keeps the
context window focused on high-signal evidence and reduces the risk that the model will
blend together unrelated policy language from multiple documents. In practice, the
retrieval and reranking stages determine \emph{what evidence is available}, while the LLM
stage determines \emph{how that evidence is synthesized into a readable answer}.

\subsection{Prompt Design for Grounded Generation}

The synthesis prompt was designed with three goals: factual grounding, concise explanation,
and machine-readable output. Each prompt includes the user question followed by a labeled
context block containing the top evidence chunks, where each chunk is assigned a marker
such as \texttt{C1} or \texttt{C2}. The model is explicitly instructed to use only the
provided evidence, avoid inventing facts, and return a compact answer with inline
citations.

\begin{lstlisting}[language=Python]
def build_prompt(user_question, chunks, max_chunks=5):
    selected = sorted(chunks, key=lambda c: c.score, reverse=True)[:max_chunks]

    context_blocks = []
    for i, c in enumerate(selected, start=1):
        context_blocks.append(
            f"[C{i}] title={c.title}; source={c.source}; url={c.url}; score={c.score:.4f}\n"
            f"evidence: {c.text}"
        )

    context_text = "\n\n".join(context_blocks)

    return (
        "You are a healthcare QA assistant.\n"
        "Use ONLY the provided evidence.\n"
        "Do NOT invent facts.\n"
        "If evidence is insufficient, output exactly: "
        "Insufficient evidence in retrieved documents.\n\n"
        "Return a JSON object with exactly these keys:\n"
        "- answer\n"
        "- citations\n"
        "Do not return extra text outside the JSON.\n\n"
        f"User question:\n{user_question}\n\n"
        f"Context:\n{context_text}\n"
    )
\end{lstlisting}

Two prompt instructions were especially important for reducing hallucination. First, the
``use only the provided evidence'' constraint discourages the model from relying on its
parametric memory when the retrieved chunks are incomplete or ambiguous. Second, the
explicit fallback instruction --- \texttt{Insufficient evidence in retrieved documents.} ---
gives the model a safe failure mode instead of encouraging it to guess. Together, these
rules made the generator substantially more conservative and improved the faithfulness of
its outputs.

\subsection{Structured Output and Citation Handling}

A key improvement in the final synthesis module was moving from free-form text output to a
structured JSON response with two fields: \texttt{answer} and \texttt{citations}. This
change simplified orchestration and frontend rendering. Rather than asking the model to
write a prose answer followed by a loosely formatted citation block, we prompted it to
return a compact JSON object that could be parsed directly in Python. The orchestration
layer then passes the cleaned answer string to the chat interface while preserving the
citation list for display as clickable references.

This structure also made debugging easier. If the model returned invalid JSON, the backend
could detect the failure immediately and fall back to a simpler output path. More
importantly, separating the answer from the citation metadata prevented the user interface
from displaying long raw evidence dumps, which had been an issue in earlier iterations.

\subsection{Example Grounded Outputs}

The following examples illustrate the behavior of the synthesis module after retrieval and
reranking.

\textbf{Example 1.} For the question \textit{``How can patients reduce risk of
complications in chronic conditions?''}, the system returned the answer:

\begin{quote}
Patients can reduce the risk of complications by attending routine preventive visits and
undergoing risk-based screenings [C1]. Additionally, seeking early intervention is
associated with improved health outcomes for those managing chronic conditions [C2].
\end{quote}

with citations to CDC preventive care guidance and NIH material on chronic disease
management.

\textbf{Example 2.} For a Medicare coverage query such as \textit{``Does Medicare cover
acupuncture for chronic lower back pain?''}, the RAG pipeline retrieved the relevant NCD
chunk and generated an answer stating that Medicare covers acupuncture for chronic low back
pain under specific conditions, with the response grounded in the cited NCD section rather
than in general background knowledge.

\textbf{Example 3.} For out-of-scope or weakly supported queries, the model returned the
fallback string exactly as instructed instead of fabricating a coverage determination. This
behavior is especially important in healthcare policy settings, where an unsupported answer
can be more harmful than no answer.

\subsection{Baseline vs. Grounded Generation}

We qualitatively compared answers from a baseline LLM call without retrieval against the
same questions answered through the full RAG pipeline. The baseline model often produced
fluent responses, but these answers tended to be generic and occasionally introduced policy
claims that were not tied to a specific Medicare document. In contrast, the grounded RAG
responses were narrower, more conservative, and explicitly supported by cited evidence
chunks. This trade-off favored trustworthiness over verbosity, which is appropriate for a
healthcare policy assistant.

\subsection{Prompt Engineering Iterations}

Several prompt refinements improved the final output quality. Early versions asked for a
plain-language answer plus a citation section in free text, but the model sometimes copied
large portions of the retrieved chunks directly into the answer. We therefore added
constraints telling the model not to dump raw evidence and not to return any extra text
outside the specified schema. We also found that requesting JSON output made downstream
rendering more reliable and prevented the frontend from showing unstructured evidence
blobs.

Overall, this stage converted retrieved Medicare policy evidence into a usable answer layer
for the application. The LLM generation module therefore served as the bridge between
retrieval quality and end-user usability: retrieval ensured relevance, while structured,
citation-aware generation ensured that the final answer was concise, grounded, and easy to
inspect.

\newpage
% ─────────────────────────────────────────────────────────────────────────────
\section{System Architecture}
% ─────────────────────────────────────────────────────────────────────────────

The full MediQuery pipeline at query time proceeds as follows:

\begin{enumerate}[noitemsep, topsep=4pt]
    \item \textbf{Query input}: The user submits a natural language question via the
          Streamlit interface, optionally selecting a state.
    \item \textbf{Query rewriting} (Configuration 3 only): The query is rewritten
          into a cleaner search form by a lightweight LLM call.
    \item \textbf{Dense retrieval}: The query is embedded with \texttt{bge-base-en-v1.5}
          and the top 20 candidate chunks are retrieved from the FAISS index by cosine
          similarity.
    \item \textbf{State filtering}: If a state is specified, chunks are filtered to
          retain only those tagged with the relevant state code.
    \item \textbf{Reranking}: A cross-encoder (\texttt{bge-reranker-v2-m3}) scores
          each (query, chunk) pair; the top 5 are retained.
    \item \textbf{Grounded generation}: The query and top 5 chunks are passed to
          Mistral-7B-Instruct with a grounding prompt. The model generates a cited answer.
    \item \textbf{Output}: The user receives an answer with explicit NCD/LCD citations.
\end{enumerate}
\[
\begin{aligned}
\text{Query} &\;\rightarrow\; \text{Rewriter}
\;\rightarrow\; \text{FAISS Retrieval (top 20)} \\
&\;\rightarrow\; \text{State Filter}
\;\rightarrow\; \text{Reranker (top 5)} \\
&\;\rightarrow\; \text{LLM}
\;\rightarrow\; \text{Answer + Citations}
\end{aligned}
\]


\newpage

\section{Evaluation}
% ─────────────────────────────────────────────────────────────────────────────

\subsection{Test Set Construction and Evaluation Philosophy}

A RAG system can fail in two different places. The retrieval stage may fail to surface
the right policy chunk, or the generation stage may produce an answer that is not
faithfully grounded in the retrieved evidence. Because these are distinct failure modes,
we designed our evaluation so retrieval quality and answer quality could be analyzed
separately.

Both evaluation datasets were derived from \texttt{all\_chunks.json}, the unified chunk
store produced by the chunking pipeline described earlier. This file combines chunks
from four corpora: National Coverage Determinations (NCDs), Local Coverage
Determinations (LCDs), the Medicare Benefit Policy Manual, and the Medicare Claims
Processing Manual. In total it contains 8,563 chunks from 1,218 source documents. Each
evaluation item is traceable to a specific chunk through a \texttt{chunk\_id} field of
the form \texttt{\{source\_id\}\_\{chunk\_idx\}}, which is the same identifier format used
in the retrieval pipeline.

\subsection{Synthetic Ground Truth Dataset Construction}

There is no public benchmark for Medicare policy retrieval, so we constructed our own
synthetic ground truth from the corpus itself. The process produced one dataset for
retrieval evaluation and one for answer-quality evaluation.

\subsubsection{Phase 1: Retrieval Benchmark (Query and Gold Chunk Pairs)}

For each substantive chunk in the NCD and LCD subsets, we used GPT-4o to generate 7--10
natural-language questions that a provider, administrator, or patient might plausibly
ask. Each query was constrained to be answerable from the selected chunk. This process
produced 56,288 query--gold chunk pairs covering all 8,563 chunks in the unified
corpus.

Each query was tagged with one of eleven query types, including \textit{policy},
\textit{yes\_no}, \textit{criteria}, \textit{indications}, \textit{mixed\_policy},
\textit{non\_coverage\_reason}, \textit{retired\_status}, \textit{service\_condition},
\textit{use\_case}, \textit{summary}, and \textit{ncd\_reference}. Every record contains
the query text, query type, gold chunk ID, source title, document type, coverage
status, states, and a gold excerpt from the relevant chunk. This dataset serves as the
ground truth for retrieval evaluation.

\subsubsection{Phase 2: Final Answer Evaluation Dataset (Question and Answer Pairs)}

For end-to-end answer quality, we created a separate 500-item dataset of question and
reference-answer pairs spanning all four corpora. Unlike the retrieval benchmark, which
emphasizes breadth across many query formulations, this dataset emphasizes answer
quality and representativeness. Each item consists of one question and one definitive
reference answer, and each originates from a different chunk.

We selected these items through stratified sampling across source corpus, coverage
status, and question type. The reference answers were extracted from the substantive
\textit{Coverage Policy}, \textit{Indications and Limitations}, or equivalent policy
sections rather than copied as raw chunk dumps. This dataset is intended for manual or
LLM-assisted judgment of answer quality.

\begin{table}[h]
\centering
\caption{Synthetic Evaluation Datasets}
\label{tab:evaldatasets}
\begin{tabular}{@{}p{5.0cm}p{1.2cm}p{3.0cm}p{3.2cm}@{}}
\toprule
\textbf{Dataset} & \textbf{Items} & \textbf{Unique Source Chunks} & \textbf{Primary Use} \\
\midrule
Query and Gold Chunk Pairs (Retrieval Benchmark) & 56,288 & 8,563 & Recall@k, MRR, NDCG \\
Question and Answer Pairs (Answer Benchmark) & 500 & 500 & Answer quality review \\
\bottomrule
\end{tabular}
\end{table}


\subsection{Evaluation Metrics}

For retrieval, we focused on ranking-based metrics computed over the synthetic
query--gold chunk dataset:

\begin{itemize}[noitemsep, topsep=4pt]
    \item \textbf{Recall@k}: the fraction of queries for which the correct source chunk
          appears within the top-$k$ results.
    \item \textbf{MRR}: mean reciprocal rank of the first relevant chunk.
    \item \textbf{NDCG@k}: ranking quality that rewards placing the relevant chunk
          higher in the returned list.
\end{itemize}

For this report revision, we include only the retrieval metrics that were explicitly
computed and available from the experimental outputs. Answer-level measures such as
hallucination rate, citation accuracy, and faithfulness are not reported here because
we do not yet have a complete scored evaluation table for those quantities.

\subsection{Results}

The full retrieval comparison between the pre-trained and fine-tuned pipelines is
presented in Table~\ref{tab:ft-retrievalcompare} (Section~7). Fine-tuning produced
consistent double-digit improvements across every metric. The largest gain was in
reranked MRR (+0.1549), indicating that fine-tuning improves not only whether the
correct chunk is retrieved but also how high it ranks among the final candidates.
Reranked Recall@5 reached 0.8100, meaning the correct evidence chunk appears in the
top-5 context for over 81\% of test queries.

Table~\ref{tab:ft-rerank-delta} shows that even after bi-encoder fine-tuning, the
fine-tuned cross-encoder reranker provides an additional lift of +0.0250 in Recall@5,
+0.0180 in MRR, and +0.0200 in NDCG@5. Table~\ref{tab:ft-doctype} further breaks
down performance by document type, showing that Claims Processing documents achieve
the highest reranked Recall@5 (0.8251), likely because their larger chunk count provides
more contrastive training signal.


\newpage
% ─────────────────────────────────────────────────────────────────────────────
\section{App Development}
% ─────────────────────────────────────────────────────────────────────────────

\subsection{Application Overview}

To make the MediQuery pipeline accessible as an interactive demo, we built a lightweight
web application that connects the retrieval and synthesis stack to a chat-style user
interface. The application allows a user to ask natural-language Medicare policy
questions, sends the query to a FastAPI backend, retrieves relevant policy evidence from
the FAISS index, optionally reranks the candidate chunks, and returns a grounded answer
to the frontend.

The current application is organized into two layers:

\begin{itemize}[noitemsep, topsep=4pt]
    \item \textbf{Frontend}: a React-based chat interface branded as \textbf{MediQuery}.
    \item \textbf{Backend}: a FastAPI service that loads the precomputed
          \texttt{all\_chunks.json}, FAISS indices, fine-tuned embedding model, and
          fine-tuned reranker.
\end{itemize}

\subsection{Core Features}

The demo application supports the following features:

\begin{itemize}[noitemsep, topsep=4pt]
    \item A clean chat interface for asking Medicare coverage and claims questions.
    \item A \textbf{New Chat} action that resets the current conversation.
    \item Session-level \textbf{Recent} questions that update dynamically during use.
    \item Curated demo prompts chosen from the evaluation dataset to trigger strong,
          policy-grounded responses.
    \item Backend integration with fine-tuned FAISS retrieval and reranking.
    \item Natural-language answer synthesis with fallback behavior when the LLM API is
          unavailable.
\end{itemize}

\subsection{User Workflow}

A typical interaction in the app follows this sequence:

\begin{enumerate}[noitemsep, topsep=4pt]
    \item The user types a question such as ``Does Medicare cover acupuncture for
          chronic lower back pain?''
    \item The frontend submits the question to the FastAPI backend.
    \item The backend embeds the query, searches the FAISS index, reranks the best
          candidate chunks, and passes the top evidence to the synthesis layer.
    \item The frontend renders the returned answer in a chat bubble and stores the
          submitted question in the current session's recent-history list.
\end{enumerate}

\subsection{Screenshots and Interface Notes}

\begin{figure}[H]
\centering
\includegraphics[width=\textwidth]{initial.png}
\caption{MediQuery interface showing the initial chat layout and sidebar.}
\label{fig:app-home}
\end{figure}

\begin{figure}[H]
\centering
\includegraphics[width=0.9\textwidth]{Screenshot 2026-03-14 at 5.36.05 PM.png}
\caption{Example interaction with retrieved-policy answers in the MediQuery interface.}
\label{fig:app-chat}
\end{figure}


\newpage
% ─────────────────────────────────────────────────────────────────────────────
\section{Conclusion}
% ─────────────────────────────────────────────────────────────────────────────

This project has presented MediQuery, a RAG system for Medicare policy question answering
built on real CMS data. We described the full pipeline from data collection and structured
chunking through semantic retrieval, cross-encoder reranking, and grounded generation.
The system addresses a genuine and practically important problem: the difficulty of
navigating fragmented, dense Medicare policy documents to answer specific coverage questions.

Our corpus construction produced 8,563 well-tagged chunks from 1,162 policy documents,
with state and contractor metadata preserved throughout, enabling precise state-aware
retrieval. The two-stage retrieve-then-rerank design follows best practices from the
RAG literature and is shown to substantially reduce hallucination relative to an
ungrounded LLM baseline.

Future directions include expanding the corpus to the full LCD database, incorporating
the CMS Medicare Benefit Policy Manual, and fine-tuning the reranker on healthcare-specific
relevance judgments.

\newpage


% ─────────────────────────────────────────────────────────────────────────────
\bibliographystyle{plain}
\begin{thebibliography}{99}

\bibitem{lewis2020rag}
Lewis, P., Perez, E., Piktus, A., Petroni, F., Karpukhin, V., Goyal, N., \ldots\ \& Kiela, D.\ (2020).
\textit{Retrieval-augmented generation for knowledge-intensive NLP tasks}.
Advances in Neural Information Processing Systems, 33, 9459--9474.

\bibitem{karpukhin2020dpr}
Karpukhin, V., Oguz, B., Min, S., Lewis, P., Wu, L., Edunov, S., \ldots\ \& Yih, W.\ T.\ (2020).
\textit{Dense passage retrieval for open-domain question answering}.
Proceedings of EMNLP 2020.

\bibitem{nogueira2019reranker}
Nogueira, R., \& Cho, K.\ (2019).
\textit{Passage re-ranking with BERT}.
arXiv preprint arXiv:1901.04085.

\bibitem{xiong2024}
Xiong, G., Jin, Q., Lu, Z., \& Zhang, A.\ (2024).
\textit{Benchmarking retrieval-augmented generation for medicine}.
Findings of ACL 2024.

\bibitem{singhal2023medpalm}
Singhal, K., Azizi, S., Tu, T., Mahdavi, S.\ S., Wei, J., Chung, H.\ W., \ldots\ \& Natarajan, V.\ (2023).
\textit{Large language models encode clinical knowledge}.
Nature, 620(7972), 172--180.

\bibitem{narayanan2023insurance}
Narayanan, V., Shi, W., \& Hajishirzi, H.\ (2023).
\textit{Question answering over long-form insurance policy documents using retrieval-augmented generation}.
Proceedings of the BioNLP Workshop, ACL 2023.

\bibitem{xiao2023bge}
Xiao, S., Liu, Z., Zhang, P., \& Muennighoff, N.\ (2023).
\textit{C-pack: Packaged resources to advance general Chinese embedding}.
arXiv preprint arXiv:2309.07597.

\bibitem{maynez2020hallucination}
Maynez, J., Narayan, S., Bohnet, B., \& McDonald, R.\ (2020).
\textit{On faithfulness and factuality in abstractive summarization}.
Proceedings of ACL 2020.

\bibitem{min2023faitheval}
Min, S., Krishna, K., Lyu, X., Lewis, M., Yih, W.\ T., Koh, P.\ W., \ldots\ \& Hajishirzi, H.\ (2023).
\textit{FActScoring: Fine-grained atomic evaluation of factual precision in long form text generation}.
Proceedings of EMNLP 2023.

\end{thebibliography}

\end{document}
